% TOP_PROBLEM: {"id": 8700026, "title": "Conjecture 3.8 — Gel’fond", "category_id": 3, "category": {"id": 3, "name": "graph_theory", "display_name": "Graph Theory", "description": "Problems involving graphs, networks, and their properties.", "slug": "graph-theory", "order_index": 3, "created_at": "2026-07-31T15:26:25.671Z"}, "status": "open", "problem_number": "AMR-086-0026", "set_id": 13, "statusQualification": "Restores the primary source setup: alpha is nonzero algebraic, beta algebraic irrational, and one fixed nonzero logarithm defines the power."}
% FIRST_POSTED: 2026-10-01
% ENTRY_KIND: statement_only
% UnsolvedMath ID 8700026; source catalogue code AMR-086-0026
% !TeX program = lualatex
% Definition and sources only; this file is not an LLM solution attempt.
\documentclass[11pt,a4paper]{article}
\usepackage[margin=25mm]{geometry}
\usepackage{fontspec}
\setmainfont{lmroman10-regular.otf}[BoldFont=lmroman10-bold.otf,
 ItalicFont=lmroman10-italic.otf,BoldItalicFont=lmroman10-bolditalic.otf]
\usepackage{amsmath,amssymb,mathtools}
\usepackage{unicode-math}
\setmathfont{latinmodern-math.otf}
\AtBeginDocument{\let\setminus\smallsetminus}
\usepackage{xurl,hyperref}
\hypersetup{colorlinks=true,urlcolor=blue}
\setcounter{secnumdepth}{0}
\setlength{\parindent}{0pt}
\setlength{\parskip}{5pt}
\setlength{\emergencystretch}{3em}
\begin{document}
\section*{Conjecture 3.8 — Gel’fond}\label{8700026}
{\small UnsolvedMath ID 8700026 · AMR-086-0026 · Graph Theory · AMR Open Problem Lists}
\subsection{Definitions and mathematical statement}
Let $\alpha\in\overline{\mathbb Q}^{\times}$ be nonzero algebraic and choose a nonzero complex logarithm $\lambda$ of $\alpha$, so that $e^\lambda=\alpha$ and $\lambda\ne0$. Let $\beta\in\overline{\mathbb Q}\setminus\mathbb Q$ be algebraic irrational. Define the chosen value of the algebraic power by
\[
\alpha^\beta:=e^{\beta\lambda}.
\]
The fixed logarithm is part of the data; a different logarithm can produce a different power.

Are $\lambda$ and $e^{\beta\lambda}$ algebraically independent over $\mathbb Q$ for every such choice? In explicit polynomial language, the conjecture states
\[
F(\lambda,e^{\beta\lambda})\ne0
\quad\text{for every nonzero }F\in\mathbb Q[X,Y].
\]
Equivalently, the field generated by these two complex numbers has transcendence degree two over $\mathbb Q$.

The algebraic-irrational hypothesis on $\beta$ and nonzero-logarithm hypothesis on $\lambda$ come from the source's surrounding setup, which the short catalogue statement omits. They are essential to the intended formulation. One does not impose an additional condition $\alpha\ne1$: the source permits $\alpha=1$ with a nonzero logarithm such as $2\pi i$.

The conclusion is a joint assertion about the logarithm and the selected power. Knowing that each entry separately is transcendental would not exclude a polynomial relation linking them. The problem quantifies over all algebraic inputs and all allowed branches, rather than a single principal-branch convention.
\subsection{Short English statement}
Are a nonzero logarithm of an algebraic number and its corresponding irrational algebraic power always algebraically independent?
\subsection{Sources}
\begin{itemize}
\item[S1] ULAM.AI and the UnsolvedMath Contributors, supplied problems.json, record 8700026 (AMR-086-0026), AMR Open Problem Lists. Catalogue identity and source transcription; CC BY 4.0.
\url{https://huggingface.co/datasets/ulamai/UnsolvedMath}.
\item[S2] Waldschmidt - Open Diophantine Problems (2004), Conjecture 3.8. Original source indexed by AMR-086-0026.
\url{https://arxiv.org/abs/math/0312440}.
\item[S3] M. Waldschmidt, Open Diophantine Problems, Moscow Mathematical Journal 4 (2004), 245–305. Author PDF supplies the surrounding definitions and hypotheses.
\url{https://webusers.imj-prg.fr/~michel.waldschmidt/articles/pdf/odp.pdf}.
\end{itemize}
\end{document}
